\makeatletter
\def\rap@fontpath{../../../Common/}
\makeatother
\documentclass{../../../Common/Latex/Templates/rapport}
\usepackage{enumitem}
\usepackage{array}
\usepackage{tabularx}
\usepackage{amssymb}

\reporttitle{Responsibility \& Signature Contract}
\projectname{Eclipse NMBU}
\orglogo{../../../Common/Eclipse Images/EclipseFullText}

% ---------------------------------------------------------------------------
% Cover / title page for rapport-class documents.
%
% Usage (after \documentclass{rapport} and before \begin{document}):
%   % ---------------------------------------------------------------------------
% Cover / title page for rapport-class documents.
%
% Usage (after \documentclass{rapport} and before \begin{document}):
%   % ---------------------------------------------------------------------------
% Cover / title page for rapport-class documents.
%
% Usage (after \documentclass{rapport} and before \begin{document}):
%   \input{../cover/rapport-cover}
%   \missionpatch{path/to/patch}      % optional circular mission/team patch;
%                                     % without it, the org logo is shown large
%                                     % at the top instead
%   \coverkind{Vedtekter}             % optional small line above the name
%   \reportsubtitle{Optional subtitle line}
%   \orgname{Eclipse NMBU}{Norwegian University of Life Sciences, Norway}
%   \reportlocation{Ås, Norge}
%   \date{3. oktober 2022}            % optional, defaults to \today
%   \contactinfo{Contact: post@eclipsenmbu.no}  % optional, shown at the bottom
%
% Then, as the first thing inside \begin{document}:
%   \makecoverpage
%
% The large title line shows \projectname (falling back to \reporttitle if
% \coverkind is not set), so a document typically sets both:
%   \projectname{Eclipse NMBU}
%   \coverkind{Vedtekter}
% ---------------------------------------------------------------------------

\makeatletter
\def\@missionpatch{}
\newcommand{\missionpatch}[1]{\def\@missionpatch{#1}}

\def\@coverkind{}
\newcommand{\coverkind}[1]{\def\@coverkind{#1}}

\def\@reportsubtitle{}
\newcommand{\reportsubtitle}[1]{\def\@reportsubtitle{#1}}

\def\@orgnameone{}
\def\@orgnametwo{}
\newcommand{\orgname}[2]{\def\@orgnameone{#1}\def\@orgnametwo{#2}}

\def\@reportlocation{}
\newcommand{\reportlocation}[1]{\def\@reportlocation{#1}}

\def\@contactinfo{}
\newcommand{\contactinfo}[1]{\def\@contactinfo{#1}}

\newcommand{\makecoverpage}{%
  \begingroup
  \thispagestyle{empty}
  \begin{center}
  \ifx\@missionpatch\@empty
    \includegraphics[width=10cm]{\@orglogo}\par
  \else
    \includegraphics[width=6.5cm]{\@missionpatch}\par
  \fi
  \vspace{3.4cm}
  \ifx\@coverkind\@empty
    {\sffamily\Huge\bfseries \@reporttitle \par}
  \else
    {\sffamily\Huge\bfseries \@coverkind \par}
    {\sffamily\Huge\bfseries \@projectname \par}
  \fi
  \ifx\@reportsubtitle\@empty\else
    \vspace{1.6cm}
    {\sffamily\bfseries \@reportsubtitle \par}
  \fi
  \vspace{2cm}
  \ifx\@missionpatch\@empty\else
    \includegraphics[height=1.3cm]{\@orglogo}\par
    \vspace{0.6cm}
  \fi
  \@orgnameone \\
  \@orgnametwo \par
  \vspace{0.6cm}
  \ifx\@reportlocation\@empty
    \@date
  \else
    \@reportlocation, \@date
  \fi
  \par
  \vfill
  \ifx\@contactinfo\@empty\else
    {\small \@contactinfo \par}
  \fi
  \end{center}
  \endgroup
  \newpage
  \pagenumbering{roman}
  \pagestyle{plain}
}
\makeatother

%   \missionpatch{path/to/patch}      % optional circular mission/team patch;
%                                     % without it, the org logo is shown large
%                                     % at the top instead
%   \coverkind{Vedtekter}             % optional small line above the name
%   \reportsubtitle{Optional subtitle line}
%   \orgname{Eclipse NMBU}{Norwegian University of Life Sciences, Norway}
%   \reportlocation{Ås, Norge}
%   \date{3. oktober 2022}            % optional, defaults to \today
%   \contactinfo{Contact: post@eclipsenmbu.no}  % optional, shown at the bottom
%
% Then, as the first thing inside \begin{document}:
%   \makecoverpage
%
% The large title line shows \projectname (falling back to \reporttitle if
% \coverkind is not set), so a document typically sets both:
%   \projectname{Eclipse NMBU}
%   \coverkind{Vedtekter}
% ---------------------------------------------------------------------------

\makeatletter
\def\@missionpatch{}
\newcommand{\missionpatch}[1]{\def\@missionpatch{#1}}

\def\@coverkind{}
\newcommand{\coverkind}[1]{\def\@coverkind{#1}}

\def\@reportsubtitle{}
\newcommand{\reportsubtitle}[1]{\def\@reportsubtitle{#1}}

\def\@orgnameone{}
\def\@orgnametwo{}
\newcommand{\orgname}[2]{\def\@orgnameone{#1}\def\@orgnametwo{#2}}

\def\@reportlocation{}
\newcommand{\reportlocation}[1]{\def\@reportlocation{#1}}

\def\@contactinfo{}
\newcommand{\contactinfo}[1]{\def\@contactinfo{#1}}

\newcommand{\makecoverpage}{%
  \begingroup
  \thispagestyle{empty}
  \begin{center}
  \ifx\@missionpatch\@empty
    \includegraphics[width=10cm]{\@orglogo}\par
  \else
    \includegraphics[width=6.5cm]{\@missionpatch}\par
  \fi
  \vspace{3.4cm}
  \ifx\@coverkind\@empty
    {\sffamily\Huge\bfseries \@reporttitle \par}
  \else
    {\sffamily\Huge\bfseries \@coverkind \par}
    {\sffamily\Huge\bfseries \@projectname \par}
  \fi
  \ifx\@reportsubtitle\@empty\else
    \vspace{1.6cm}
    {\sffamily\bfseries \@reportsubtitle \par}
  \fi
  \vspace{2cm}
  \ifx\@missionpatch\@empty\else
    \includegraphics[height=1.3cm]{\@orglogo}\par
    \vspace{0.6cm}
  \fi
  \@orgnameone \\
  \@orgnametwo \par
  \vspace{0.6cm}
  \ifx\@reportlocation\@empty
    \@date
  \else
    \@reportlocation, \@date
  \fi
  \par
  \vfill
  \ifx\@contactinfo\@empty\else
    {\small \@contactinfo \par}
  \fi
  \end{center}
  \endgroup
  \newpage
  \pagenumbering{roman}
  \pagestyle{plain}
}
\makeatother

%   \missionpatch{path/to/patch}      % optional circular mission/team patch;
%                                     % without it, the org logo is shown large
%                                     % at the top instead
%   \coverkind{Vedtekter}             % optional small line above the name
%   \reportsubtitle{Optional subtitle line}
%   \orgname{Eclipse NMBU}{Norwegian University of Life Sciences, Norway}
%   \reportlocation{Ås, Norge}
%   \date{3. oktober 2022}            % optional, defaults to \today
%   \contactinfo{Contact: post@eclipsenmbu.no}  % optional, shown at the bottom
%
% Then, as the first thing inside \begin{document}:
%   \makecoverpage
%
% The large title line shows \projectname (falling back to \reporttitle if
% \coverkind is not set), so a document typically sets both:
%   \projectname{Eclipse NMBU}
%   \coverkind{Vedtekter}
% ---------------------------------------------------------------------------

\makeatletter
\def\@missionpatch{}
\newcommand{\missionpatch}[1]{\def\@missionpatch{#1}}

\def\@coverkind{}
\newcommand{\coverkind}[1]{\def\@coverkind{#1}}

\def\@reportsubtitle{}
\newcommand{\reportsubtitle}[1]{\def\@reportsubtitle{#1}}

\def\@orgnameone{}
\def\@orgnametwo{}
\newcommand{\orgname}[2]{\def\@orgnameone{#1}\def\@orgnametwo{#2}}

\def\@reportlocation{}
\newcommand{\reportlocation}[1]{\def\@reportlocation{#1}}

\def\@contactinfo{}
\newcommand{\contactinfo}[1]{\def\@contactinfo{#1}}

\newcommand{\makecoverpage}{%
  \begingroup
  \thispagestyle{empty}
  \begin{center}
  \ifx\@missionpatch\@empty
    \includegraphics[width=10cm]{\@orglogo}\par
  \else
    \includegraphics[width=6.5cm]{\@missionpatch}\par
  \fi
  \vspace{3.4cm}
  \ifx\@coverkind\@empty
    {\sffamily\Huge\bfseries \@reporttitle \par}
  \else
    {\sffamily\Huge\bfseries \@coverkind \par}
    {\sffamily\Huge\bfseries \@projectname \par}
  \fi
  \ifx\@reportsubtitle\@empty\else
    \vspace{1.6cm}
    {\sffamily\bfseries \@reportsubtitle \par}
  \fi
  \vspace{2cm}
  \ifx\@missionpatch\@empty\else
    \includegraphics[height=1.3cm]{\@orglogo}\par
    \vspace{0.6cm}
  \fi
  \@orgnameone \\
  \@orgnametwo \par
  \vspace{0.6cm}
  \ifx\@reportlocation\@empty
    \@date
  \else
    \@reportlocation, \@date
  \fi
  \par
  \vfill
  \ifx\@contactinfo\@empty\else
    {\small \@contactinfo \par}
  \fi
  \end{center}
  \endgroup
  \newpage
  \pagenumbering{roman}
  \pagestyle{plain}
}
\makeatother

\missionpatch{../../../Common/Eclipse Images/EclipseMoon}
\coverkind{Responsibility \& Signature Contract (ansvarsog signaturkontrakt)}
\reportsubtitle{New Bylaws version: ECL-H2026-RESP-01B}
\orgname{Eclipse NMBU}{Norwegian University of Life Sciences, Norway}
\reportlocation{Ås}
\date{\today}
\contactinfo{Contact: post@eclipsenmbu.no}

\begin{document}

\makecoverpage
\tableofcontents
\reportmainmatter

\noindent\textbf{Document ID:} ECL-H2026-RESP-01B\\
\textbf{Applies from:} the New Bylaws' entry into force (proposed New Bylaws \S17.2). If the Bylaw Amendment fails, use ECL-H2026-RESP-01A instead.\\
\textbf{Governing Bylaws:} proposed New Bylaws \S14 (Mandate, delegation and Signature Right), \S4.3 no.\ 7--8 (Nearest Leader, Responsibility Role), \S5.1 (Core Member).

\vsection{1. Core principle}
Signing a document on behalf of Eclipse NMBU is itself a responsibility. Signature Right is therefore not merely an administrative permission: it is assigned through this same Responsibility system, alongside any other genuine leadership responsibility (Master Drafting Prompt \S17). Responsibility, delegation, Authority and Signature Right are treated as one responsibility relationship, since the proposed New Bylaws do not distinguish them into separate regimes.

\vsection{2. Parties}
\begin{tabularx}{\textwidth}{|>{\bfseries\raggedright\arraybackslash}p{3cm}|X|}
  \hline
  From & \rule{0pt}{1cm}The person currently holding the responsibility (or, where the General Assembly itself is the granting party, see \S7 below). \\
  \hline
  To   & \rule{0pt}{1cm}The person receiving the responsibility.                                                                                      \\
  \hline
\end{tabularx}

The transfer is signed by both parties, then logged by the Board (see Board Record, ECL-H2026-RESP-02). \textbf{The transfer is effective automatically once properly signed by both parties}: the Board's role under this contract is to record the transfer, not to approve it. \emph{[NOTE: neither Bylaws version gives the Board a veto over a properly-signed transfer between two members; this is the resolved, Board-confirmed position for this document set: see the Clarifications Needed note, ECL-H2026-README. The Board retains its ordinary powers under \S5.7/\S14.5 (e.g. to change who holds an office, or to suspend a delegated Mandate by 2/3 vote) independently of this contract.]}

\vsection{3. Responsibility direction}
Responsibility may only move \textbf{downward} through the organisational hierarchy (From a leader To someone reporting to them, directly or via the chain in proposed New Bylaws \S4.3 no.\ 7). It cannot move sideways between peers. It can move upward only through the Escalation mechanism in \S6 below, when the responsible superior is unavailable: not as an ordinary transfer.

\vsection{4. Responsibility being transferred}
\begin{tabularx}{\textwidth}{|>{\bfseries\raggedright\arraybackslash}p{4.5cm}|X|}
  \hline
  Scope                                & \rule{0pt}{1.6cm}                                                                                                                                             \\
  \hline
  Specific responsibilities            & \rule{0pt}{1.8cm}                                                                                                                                             \\
  \hline
  Expected outcomes                    & \rule{0pt}{1.6cm}                                                                                                                                             \\
  \hline
  Delivery obligation (leveringsplikt) & \rule{0pt}{1.2cm}The recipient is obliged to deliver within the set time frame (proposed New Bylaws \S5.6.2).                                                 \\
  \hline
  Reporting obligation (meldeplikt)    & \rule{0pt}{1.2cm}If the recipient cannot carry out a task taken on, they shall notify their Nearest Leader without undue delay (proposed New Bylaws \S5.6.3). \\
  \hline
  Instance-specific upward obligations & \rule{0pt}{1.6cm}Any additional reporting or approval steps specific to this responsibility (e.g. financial limits, safety sign-off): list here.            \\
  \hline
  Signature Right included?            & \rule{0pt}{1cm}Yes \quad $\square$ \quad\quad No \quad $\square$ \quad\quad If yes: within what limits? \rule{3cm}{0.4pt} (proposed New Bylaws \S14.9)        \\
  \hline
\end{tabularx}

\vsection{5. Obligations of the person giving responsibility (From)}
The recipient (To) can expect the following from the person giving responsibility (From):
\begin{enumerate}
  \item Clear intent: what is actually being handed over, and why.
  \item Real delegated Authority: not just a task, but the standing to make the decisions the responsibility requires.
  \item The resources/access required to carry out the responsibility.
  \item Availability from the From party for questions and support.
  \item Support where the recipient's own responsible team falls short.
  \item Honest feedback on performance in the role.
\end{enumerate}

\vsection{6. Escalation}
If the person marked \textbf{From} is unavailable, escalation goes to that person's \textbf{Nearest Leader (nærmeste leder)}, per proposed New Bylaws \S4.3 no.\ 7 (member \(\to\) Team Lead \(\to\) Project Lead \(\to\) responsible Board member). \emph{No fixed Deputy Chair or CTO escalation rule applies here} except where the CTO is in fact the Nearest Leader in that specific chain, or is acting as Deputy Chair for the CEO under \S12.3.

\vsection{7. Core Membership}
Receiving a genuine leadership responsibility through this contract makes the recipient a \textbf{Core Member (kjernemedlem)} (proposed New Bylaws \S5.1 no.\ 2). Ordinary members do not receive this tier merely by being asked to complete a task. \textbf{Signing an ordinary task/order is not, by itself, a Responsibility Contract}: this contract is used only where a genuine, standing leadership responsibility (e.g. Project Lead, Team Lead, or delegated Signature Right) is being transferred, not for one-off task assignments.

\vsection{8. General Assembly \(\to\) CEO}
Where the General Assembly itself is the granting party (e.g. the Mandate the newly elected CEO holds under \S14.1), the contract is worded:

\begin{quote}
  ``The General Assembly, per the protocol of [DATE], ECL-H2026-PROT-01, gives [CEO NAME] the responsibility of [SCOPE].''
\end{quote}

The CEO signs as \textbf{To}. The Protocol Signers who certified the General Assembly's decision \textbf{do not} personally become the \textbf{From} party: the responsibility originates from the General Assembly as a whole, and the Protocol ID above is the record of that origin.

\vsection{9. Board record}
Once signed, this contract is logged by the Board using the Board Record template, ECL-H2026-RESP-02. This contract and the Board Record cross-reference each other by Document ID.

\vsection{10. Signatures}
\begin{tabularx}{\textwidth}{|>{\bfseries\raggedright\arraybackslash}p{4cm}|X|}
  \hline
  From (name, role)        & \rule{0pt}{1.2cm} \\
  \hline
  From signature and date  & \rule{0pt}{1.2cm} \\
  \hline
  To (name, role)          & \rule{0pt}{1.2cm} \\
  \hline
  To signature and date    & \rule{0pt}{1.2cm} \\
  \hline
  Board Record Document ID & \rule{0pt}{1cm}   \\
  \hline
\end{tabularx}

\end{document}
